(1) 记 $r_+=\sqrt{x^2+y^2+(z-d)^2}$、$r_-=\sqrt{x^2+y^2+(z+d)^2}$。在上介质可取 $\Phi_1=(q/r_++q'/r_-)/(4\pi\varepsilon_1)$；下半空间用无源传递等效势 $\Phi_2=q''/(4\pi\varepsilon_2r_+)$，因为源点在下半空间域外。

(2) 界面 $z=0$ 上 $r_+=r_-$，电势连续给 $(q+q')/\varepsilon_1=q''/\varepsilon_2$。法向位移连续要求 $q-q'=q''$，故 $\boxed{q'=q(\varepsilon_1-\varepsilon_2)/(\varepsilon_1+\varepsilon_2)}$，$\boxed{q''=2\varepsilon_2q/(\varepsilon_1+\varepsilon_2)}$，两侧边值条件同时满足。

(3) 真实电荷在位置 $z=d$ 只受上域镜像等效电荷的附加场，$\boxed{F_z=qq'/(16\pi\varepsilon_1d^2)}$。当 $\varepsilon_2>\varepsilon_1$ 时 $q'<0$，力向负 $z$、吸引界面；若 $\varepsilon_2<\varepsilon_1$ 则镜像同号，粒子受排斥。

(4) 界面诱导势在电荷处为 $\Phi_{\rm ind}=q'/(8\pi\varepsilon_1d)$。正确相互作用能是 $U=\frac12q\Phi_{\rm ind}=qq'/(16\pi\varepsilon_1d)$，故 $-dU/dd=qq'/(16\pi\varepsilon_1d^2)=F_z$。$\varepsilon_2=\varepsilon_1$ 时无界面力；$\varepsilon_2\to\infty$ 时 $q'\to-q$，恢复接地理想导体平面镜像结果。
